% Source: https://artofproblemsolving.com/community/c3h1582862_17_geometry_questions
% Stable identity: aops-topic-1582862-q12
One lap consists of the two straightaways and two semicircles, which together make one full circle of diameter $40$ m:

\begin{asy}
usepackage("amsmath");
size(6cm);
defaultpen(fontsize(10pt));
real r=2;
draw((0,-r)--(5,-r),linewidth(0.9));draw((0,r)--(5,r),linewidth(0.9));
draw(arc((5,0),r,-90,90),linewidth(0.9));draw(arc((0,0),r,90,270),linewidth(0.9));
draw((5,-r)--(5,r),dashed);draw((0,-r)--(0,r),dashed);
label("$100$ m",(2.5,r),N);label("$100$ m",(2.5,-r),S);label("$40$ m",(5,0),W);label("$40$ m",(0,0),E);
label("\mbox{}",(min(currentpicture,user=true)+max(currentpicture,user=true))/2,black); // keeps the text after the figure black when LaTeX typesets the labels inline
\end{asy}

\[
\text{lap}=2\cdot100+40\pi\approx200+125.66=325.66\text{ m}.
\]
Four laps are about $1302.65$ m, so the average speed is
\[
\dfrac{4(200+40\pi)}{6}\approx\dfrac{1302.65}{6}\approx217.11\text{ m per minute}\approx\boxed{217\text{ m per minute}}.
\]
